\uuid{3BvB}
\titre{Trigonométrie}
\auteur{bodin ; hanani ; mzari}
\organisation{exo7}
\licence{CC BY-NC-SA 4.0}
\datecreate{2020-01-20}
\niveau{L1}
\module{Algèbre}
\chapitre{Nombres complexes}
\sousChapitre{Forme cartésienne, forme polaire}
\difficulte{2}
\theme{Écritures algébrique et géométrique, trigonométrie}
\exo7qcmid{lille-1.nombres-complexes.010}

\contenu{
\begin{question}
\qclassification{104.01}
\qtags{motcle=trigonométrie}
Soit $\theta$ un réel.  Quelles sont les assertions vraies ?
\begin{answers}

     \bad{$\cos^3\theta= \frac{1}{8}(\cos(3\theta) +3\cos \theta)$}

    \good{$\cos^3\theta= \frac{1}{4}(\cos(3\theta) + 3\cos \theta)$}

     \good{$\sin^3\theta= \frac{1}{4}(3\sin \theta - \sin(3\theta))$}

      \bad{$\sin^3\theta= \frac{1}{4}(3\sin \theta + \sin(3\theta))$}

\end{answers}
\begin{explanations}
On peut appliquer les formules d'Euler.

\end{explanations}
\end{question}
}
